% Slide — SA antes/depois
\begin{frame}{Simulated Annealing -- \antes\ vs. \depois}
  \begin{columns}[T,onlytextwidth]
    \column{0.48\textwidth}
      \antes\\[0.3em]
      \small \texttt{scipy.optimize.dual\_annealing}
      \begin{itemize}\footnotesize
        \item Cronograma de temperatura e critério de aceitação
          \textbf{internos à biblioteca} -- caixa-preta com parâmetros
          próprios, não os documentados na proposta.
        \item Refinamento local por Nelder-Mead a cada iteração.
        \item Uma cadeia, avaliação sequencial.
      \end{itemize}
    \column{0.48\textwidth}
      \depois\\[0.3em]
      \small \texttt{TorchSA}: cronograma explícito
      \begin{itemize}\footnotesize
        \item $T_0=1000$, $\alpha_{\text{resfr.}}=0{,}95$ (Seção 4.4 da
          proposta), $T_{k+1}=\alpha\, T_k$.
        \item Critério de Metropolis auditável: aceita se $\Delta C<0$, ou
          com probabilidade $e^{-\Delta C/T_k}$.
        \item \highlight{32 cadeias independentes} (\texttt{n\_chains=32})
          em paralelo, mesmo cronograma, saída = melhor entre elas.
      \end{itemize}
  \end{columns}
  \vspace{0.8em}

  \begin{block}{Por que isso é uma reescrita, não só uma correção}
    \small SA é inerentemente sequencial. Vetorizar sem descaracterizar o
    método significa rodar várias cadeias \textbf{independentes} com o
    mesmo cronograma, em vez de fingir paralelismo dentro de uma cadeia.
    O ganho de robustez vem de reportar o melhor entre 32 buscas, não de
    uma cadeia mais rápida.
  \end{block}

  \note{
    O scipy.dual\_annealing é uma implementação respeitável, mas é uma
    caixa-preta: o cronograma de temperatura e os parâmetros de aceitação
    internos não são os mesmos documentados na Seção 4.4 da proposta, o
    que dificultava justificar por que aquele T0 e aquele alpha específicos.
    A versão atual implementa o critério de Metropolis manualmente, o que
    o torna auditável linha a linha, e explora paralelismo de um jeito
    metodologicamente honesto: como SA clássico é sequencial por natureza,
    a vetorização aqui significa rodar 32 cadeias independentes com o
    MESMO cronograma de temperatura, não inventar um SA "paralelo" que não
    existe na literatura -- a saída é simplesmente a melhor solução entre
    as 32.
  }
\end{frame}
